
% Abstract

Increased reliance by AI systems in high-risk public sector domains has precipitated the emergence of data pipelines whose architectural properties directly correlate with the dimensions of fairness, privacy, and accountability. Though current data lakehouse and ETL architectures offer excellent guarantees in terms of quality and scaling properties, they fail to provide a means for thwarting the transmission of the negative effects of historical biases, privacy breaches, and legal infractions to downstream models. To fill the gap in the current state-of-the-art architectures, the FAIR-CARE Lakehouse is presented in this section, an architectural design incorporating the FAIR data principles and addressing causal validity, privacy, and fairness through the constraints of the FAIR-spectrum and the Medallion (Bronze-Silver-Gold) lifecycle. The architectural design incorporates PII identification, privacy-preserving operations, causal graph learning, and fairness constraint in feature governance in an end-to-end design pipeline facilitated by a quantifiable FAIR-CARE score measuring architectural readiness. Validation and testing of the designed architecture are conducted through the application of benchmark datasets in the domains of COMPAS, Adult, and German credit datasets and a high-risk domain study based on the challenge contest "NIJ Recidivism Competition." The results verified through experimentation and deduction suggest a decrease in the value of Equalized Odds disparities of up to 60% while satisfactorily supporting competitive predictive model capability along the lines of stringent privacy provisioning.This study contributes to the domain by developing a reusable architectural design and design scoring methodology and empirical investigation illustrating how software architects are able to direct risks associated with the socio-technical implementation and execution processes of AI-enabled strategic systems in a constant and effective manner through the disciplines of software architecture and design science theory and methodology application in developing software architectural solutions to high-stake issues in AI-enabled strategic systems associated with privacy and accountability in strategic decision-making processes and tasks explored in the domain of software architecture and AI and analytics applications associated with enabling strategic decision-making and management.

% Introduction
The growing integration of AI-driven decision-making into public-sector and safety-critical environments has introduced new pressures on software architecture. Traditional architectural concerns, performance, scalability, interoperability, are no longer sufficient when system outputs directly influence individual rights, access to resources, or legal determinations. In domains such as corrections, public safety, and healthcare, the software architect is increasingly responsible for ensuring that data infrastructures uphold principles of fairness, privacy, causal validity, and regulatory compliance. All these forces redefine the agenda of architecture, where architecture is no longer required to focus marginally better due to its technical capabilities but aims to navigate the data life cycle while keeping in mind the risks associated with socio-technical problems. To revisit and redesign architectures in a reliable and transparent manner requires well-structured and end-to-end data management pipelines, similar to the flow of data illustrated in Figure . Since AI-based systems work in a critical area where architecture should aim to be reliable along with other parameters such as robustness, transparency, and accountability and remove any systemic biases and historical biases in addition to its performance  , empirical studies have noted that algorithmic systems such as COMPAS tend to perform in a biased manner and tend to produce biases even when the protected attributes are removed  . Biased datasets in the medical domain tend to increase disparity in both the detection and treatment of diseases  . For instance, a commercial tool supporting the care and treatment of patients in the United States tends to consider the cost of treatment in place of medical need and tends to treat Black patients undertreated  . The most important takeaway from the entire discussion is that focusing on the entire architecture of health-based machine learning requires the need to take into consideration the entire life cycle  .

% Motivation: Architecting for High-Risk AI Systems
Where AI decisions can shape social, legal, or safety outcomes, the technology is not an ancillary piece, but a core component in an end-to-end system that captures and transforms information and feeds it to other downstream models for use. So, in real-risk contexts, making these systems reliable requires more than just running traditional ETL data pipelines. This means explicitly addressing non-functional requirements such as fairness, privacy, provenance, and auditability. It’s in remediation datasets that we can see why high-risk management is so difficult. Models that predict recidivism can suffer from data bias, where patterns reflect past surveillance rather than real risk. There are also huge privacy concerns: once you incorporate pseudo-identifiers like release dates, locations, and crime categories, you can re-identify individuals even after direct identifiers have been removed. Both of these problems strongly suggest that ethical and regulatory rules should be built into the architecture itself, rather than left to model developers or compliance teams in isolation.

% Problem Statement
The classical ETL and Lakehouse architectures are designed to handle fast processing and handle-schema-evolution issues well, but they tend to struggle in incorporating ethics guidelines across the whole data life cycle. The issues are clear in a number of ways:\\
- Ethical constraints, such as mitigation of bias, anonymization of data, and regulatory compliance, are distributed across various teams and do not have a unifying architectural perspective regarding the end-to-end process .\\
- Traditionally, pipelines tend to consider correlation enough to trigger the next task but fail to provide guardrails to identify and prevent misleading causal inferences, such as Zip Code → Race → Risk .\\
- Privacy and regulatory standards, for instance, GDPR Anonymization, HIPAA Expert Determination, and CCPA Functional Separation, are seldom incorporated within the pipeline logic in the DL algorithm itself, and the verification of compliance is conducted post-processing .

As a result, there is a significant need for an architecture that raises ethics, privacy, lineage, and causal reasoning to a premier status.
The way ahead involves a renaming of the software architect’s function: to make the software architect the keeper of data ethics and the guardian of a responsive and ethically informed system.

% Research Questions
This paper explores how to stitch ethical governance into data pipelines that back high-stakes AI applications. We address the following key questions:

RQ1: How can we design a complete end-to-end data pipeline within the Medallion lakehouse framework that will enforce ethical governance, fairness, privacy, and regulatory compliance at every step?

RQ2: What technical components are required at each layer-causal analysis, differential privacy, synthetic data generation and bias mitigation-towards fulfilling the requirements of CARE without an undue sacrifice in utility?

RQ3: How can we evaluate the ethical readiness of datasets within high-risk domains through the fusion of the FAIR data principles with explicit ethical criteria into one layer-aware metric?

% Contributions
To address these challenges, this paper introduces the FAIR-CARE Lakehouse, a reference architecture that embeds FAIR data principles alongside the CARE ethical framework (Causality, Anonymity, Regulatory compliance, Ethics) across the Medallion lifecycle. Our contributions are as follows:

- Reference Architecture: we propose the extension of Medallion Lakehouse to Bronze-Silver-Gold by considering PETs, such as causal inference mechanisms and clearly defined governance policies, as integral architectural components.

- Algorithmic Framework: The approach represents a multistage pipeline that deals with the processing of PII detection, DP, SDG, causal structure learning, and human-in-the-loop validation while constructing modular and reusable components that can easily be plugged into the architecture.

- Evaluation Framework: We introduce the FAIR-CARE Score, a composite, layer-specific metric that evaluates privacy leakage, fairness disparities, lineage completeness, and causal soundness to quantify ethical readiness.

- Regulatory Compliance Mapping: We demonstrate how architectural parameters (e.g., DP privacy budget $\epsilon$, k-anonymity thresholds, lineage granularity) can be configured to align with GDPR, HIPAA, and CCPA requirements.

- Empirical Validation: Using benchmark datasets from the NIJ Recidivism Forecasting Challenge  and widely studied fairness datasets, we illustrate how legacy correctional data can be transformed into a compliant, fair, and analytically valuable asset.

% Related Work

Research on trustworthy AI, privacy-preserving data engineering, and data lakehouse architectures has advanced substantially, yet existing approaches remain siloed and insufficient for architecting high-risk AI systems. Prior work in AI ethics has thoroughly documented how biased datasets, opaque transformations, and ungoverned model pipelines can propagate systemic harms . The analyses of recidivism risk tools such as COMPAS show that removing the protected attributes does not stop biased outcomes if upstream data steps embed structural hints about race or socioeconomic status. Yet, most works focus on fairness at the model level and provide little insight into how choices baked into the data pipelines fuel bias as data flows. For that reason, there are no strong pipeline-level mechanisms to keep fairness, causal validity, and transparency intact throughout the whole data lifecycle. Regulatory standards like GDPR, HIPAA, and CCPA have detailed the rules for anonymization, pseudonymization, and functional separation .

Engineers in the current settings commonly make use of ad-hoc masking, auditing by design, and strict access controls in order to satisfy the above constraints. The above methods are fragile and cumbersome and work poorly in systems with the characteristic continuous flow of data in modern architectures . The above methods mostly include last-minute testing to ascertain a compliant flow chain, opening up the potential for leakage and misuse issues while the data has passed through various systems and ends up somewhere else along the way. Current studies in ML ethics provide clear conceptual foundations, but current architectures fail to provide any means to design and implement the privacy assurances in a manner that reflects a system-wide perspective in architectural fashion and timetabling .  Studies focusing on the intersection and interconnectivity of privacy engineering in data and algorithmic fairness establish a fundational architectural toolset consisting of tools such as the application of Differential Privacy notions for privacy engineering, the incorporation of artificial/synthetic data generation, and the aim to measure Demographic Parity and Equalized Odds in metrics and studies . Causal ideas and applications, and most notably Counterfactual Fairness, present a clearly effective application and inspection tool for the discovery of spurious relations in datasets through exploratory and prototypic applications . The application and implementation of the above tools in model prototyping and applications mostly occur in isolated applications in the design and prototyping phase, and less in the implementation architecture phase in their application and methodology . The implementation of the pipeline in the architecture mostly neglects the lineage and governance application and constraints in the pipeline implementation architecture, and mostly fails to attain the capability to provide the uniform and constant implementation of the tools and methods in a valid and effective application and implementation process and solution. The current implementation of AI/ML mostly views the application of ethics and methodology concerning the implementation in the architecture and prototyping phase, mostly in the preprocessing phase in their implementation and application. The current architecture of the lakehouse implementation (Medallion architecture or the Bronze, Silver, and Gold model)  provides an effective and promising architectural implementation in the lakehouse architecture framework concerning the provision of a robust implementation concerning the capability in the evolution of the schema in the architecture and the application concerning the provision in the architecture concerning the capability in the ETL/ELT and analytic readiness in the architecture and the current implementation and availability concerning the lakehouse architecture implementation and framework. Nonetheless, the current implementation and application and availability in the architecture and the implementation mostly fail to provide a robust implementation concerning the provision in the architectural implementation and availability concerning the embedding and application and enforcement of ethical governance principles, such as differential privacy guarantees, causal analysis safeguards, synthetic data validation, and continuous fairness auditing, directly into the operational layers of the lakehouse pipeline. As a result, the Medallion architecture, despite its strengths in schema evolution, scalability, and analytic readiness, still lacks a unified, system-wide mechanism for operationalizing ethical constraints end-to-end. This gap underscores the need for an architectural reframing in which fairness, privacy, lineage, and causal reasoning become intrinsic, enforceable components of the data pipeline rather than optional or peripheral add-ons evaluated only during prototyping or post-hoc compliance checks.
The limitations of existing approaches are summarized in Table . Across domains, the literature provides conceptual tools (fairness metrics, privacy mechanisms, causal reasoning), but there is no unified architecture that:
(1) treats fairness and privacy as first-class architectural concerns;
(2) integrates PETs, causal validation, and fairness filtering directly into Medallion-layer transformations;
(3) ensures continuous, lineage-aware governance; and
(4) provides a measurable governance metric for dataset readiness.
The FAIR-CARE Lakehouse directly addresses these gaps by operationalizing ethical constraints as intrinsic pipeline properties rather than external evaluations.

In adjacent domains such as healthcare, large-scale studies highlight how ungoverned data pipelines can encode and magnify structural inequities. Obermeyer et al.  revealed pervasive racial bias in a commercial health risk model affecting over 200 million patients, rooted not in model algorithms but in upstream data proxies. Rajkomar et al.  emphasize the need for continuous lifecycle monitoring and stewardship principles parallel to those required for correctional systems. These findings further reinforce that high-risk domains demand architectural, not merely algorithmic, interventions.

Taken together, existing research provides essential conceptual components fairness metrics, differential privacy, causal reasoning, and scalable lakehouse infrastructures, but does not unify them into a governed architecture capable of supporting high-risk AI systems. The FAIR-CARE Lakehouse fills this gap by embedding ethical constraints into the structural layers of the Medallion architecture, enabling continuous, auditable, and legally aligned data governance.

% Methodology

% Architectural Overview
We propose the FAIR-CARE Lakehouse, a reference architecture that integrates FAIR data principles with the CARE ethical framework (Causality, Anonymity, Regulatory compliance, Ethics). As illustrated in Figure , the system extends the standard Medallion pattern (Bronze--Silver--Gold)  by embedding privacy-enhancing technologies (PETs) and causal inference engines directly into the transformation logic.

The architecture is implemented using a hybrid distributed computing mechanism. While Apache Spark is utilized for the more complex ETL and schema application tasks in the Bronze tier, we use modular Python components in the Silver and Gold tiers. This division allows for the integration of specialized libraries needed for causal discovery algorithms and synthetic data generation.

% Layered Design and Transformations

% Bronze Layer: The Quarantine Zone
The Bronze layer acts as the immutable system of record. Data is entered via Spark Structured Streaming, with ACID (Atomicity, Consistency, Isolation, and Durability) compliance through Delta Lake.

- PII Detection: An automated scanner using Microsoft Presidio and Regex patterns checks columns for Personally Identifiable Information (PII) and quasi-identifiers (e.g., Zip Code, Release Date).
- Pseudonymization: Direct identifiers get hashed with a salted key stored in a hardware security module (HSM). This implements the "Functional Separation" that CCPA requires .
- Provenance: Metadata about source, ingestion timestamp, and schema hash goes to the Unity Catalog.

% Silver Layer: The Privacy and Causal Engine
The Silver layer converts risky and pseudonymized data into investigation-ready assets through isolated Python tasks.

- Differential Privacy (DP): $\epsilon$-Differential Privacy is applied through the Laplace mechanism for Numerical Perturbation. The Control Plane maintains a privacy budget ($\epsilon$). Once the budget is exhausted, the dataset would be immutable to prevent reconstruction attacks .
- Adaptive Anonymization: For non-numeric attributes, the methodology is to select between $k$-anonymity, $l$-diversity, and $t$-closeness depending on the sensitivity level of the data and the distribution of the quasi-identifiers.
- Causal Discovery: Unlike traditional ETL, the work conducted within this layer builds a Causal Directed Acyclic Graph (DAG) employing algorithms like as PC or NOTEARS through CausalNex . Based on this graph, we identify the confounding variables (like Race $\rightarrow$ SES $\rightarrow$ Recidivism) so that we can correctly avoid correlations between models.

% Gold Layer: Fairness and Feature Store
The Gold layer produces the final feature store.

- Markov Blanket Selection: We filter features based on the causal DAG. Only features in the target variable's Markov Blanket make it through, ensuring causal sufficiency .
- Bias Mitigation: Pre-processing algorithms from AIF360, as in Reweighing or Disparate Impact Remover, correct historical imbalances .
- Vector Embedding: For Retrieval-Augmented Generation (RAG) applications, text fields (case narratives, for example) get vectorized and stored in Qdrant, tagged with their fairness certification.

% The FAIR-CARE Score
We define a composite score, calculated at runtime, to quantify ethical readiness. It's a weighted sum of layer-specific metrics:

\begin{equation}
\text{Score}_{FC} = w_B S_B + w_S S_S + w_G S_G
\end{equation}

The sub-scores break down as follows:

- Bronze Score ($S_B$): Measures ingestion quality.
$$S_B = \frac{1}{3}(\text{Provenance}_{%} + \text{PII}_{\text{Recall}} + \text{Quality}_{\text{Base}})$$

- Silver Score ($S_S$): Measures privacy-utility balance.
$$S_S = \frac{1}{4}(\text{Anon}_{\text{Str}} + \text{Util}_{\text{Ret}} + \text{Causal}_{\text{Valid}} + \text{HITL}_{\text{Appr}})$$
$\text{Anon}_{\text{Str}}$ comes from achieved privacy loss ($1 - \frac{\epsilon}{\epsilon_{max}}$) and $k$-anonymity thresholds.

- Gold Score ($S_G$): Measures downstream fairness.
$$S_G = \frac{1}{3}(\text{Fairness}_{\text{Met}} + \text{Feat}_{\text{Qual}} + \text{Util}_{\text{Pred}})$$
$\text{Fairness}_{\text{Met}}$ aggregates pass/fail rates for Demographic Parity and Equalized Odds.

Algorithm  shows the ``Privacy Transformer'' logic, a key Silver layer component that enforces these constraints.

% Architectural Assumptions and Quality Attribute Trade-offs
The FAIR-CARE Lakehouse architecture is based on a number of architectural assumptions that guide its design choices and quality attribute trade-offs. Primarily, the architecture assumes that privacy, fairness, and utility cannot be optimized simultaneously but instead form a tripartite trade-off space similar to the traditional “Security–Performance” trade-off space in software architecture studies. By incorporating Differential Privacy and causal filtering in the Silver layer, the system reflects a design choice that favors privacy and fairness over purely predictive utilities. This trade-off reflects a design criterion more in line with the demands of high-risk domains, wherein stringent ethics and laws take priority over the pursuit of prediction accuracy. For architectural assessment in terms of the quality attributes discussed, the pipeline architecture can be analyzed through scenario-based reasoning derived from the application of the software architecture techniques of ATAM and QAW-lite. The points of concern in the architecture include: (i) gradual loss of privacy in repeated querying, (ii) drift in fairness due to changes in the upstream feature space through software updates, and finally, (iii) loss of utility due to an increase in the values of the anonymity parameters, to name a few.
Even if a full treatment of the architecture through the lens of software architectural techniques such as ATAM is deferred in this work, the architectural design patiently argues its constituent design rationales in terms of the prioritized quality attributes. Another assumption implicitly made is that the causal structure will be validated by human experts. Though the Silver layer has the capability to automatically find the causal structure, it is the domain experts in the form of criminal-legal scholars and social scientists who should analyze and refine the DAGs produced by the Silver layer. The Human-in-the-Loop process is much more than a checkbox process and an integral architectural aspect of the pipeline instead. The pipeline will need to be halted at several crucial points to seek the validation of the domain experts before processing the data from the Silver layer to the Gold layer phase. Though this has additional costs involved, it improves the causal validity and regulatory defensibility of the process instead. All the above assumptions and trade-offs specify what the FAIR-CARE Lakehouse really means, a system that seeks no less than predictive sharpness but at a cost and a data architecture governed by ethics and causal constraints designed into it instead.

% Experimental Evaluation

% Datasets
We validate the FAIR-CARE Lakehouse architecture using four benchmark datasets in fairness and high-risk decision-making research . These datasets span correctional forecasting, credit scoring, and income prediction, domains in which fairness, privacy, and causal integrity are critical:

- NIJ Recidivism Challenge: A correctional dataset containing over 25{,}000 records with sensitive behavioral and supervision attributes, such as ground-truth recidivism outcomes .
- COMPAS: A benchmark for evaluating racial algorithmic risk assessment .
- UCI Adult \& German Credit: Standard fairness benchmarks often employed to investigate the socioeconomic bias of income and credit scoring.

These datasets enable a representative assessment across multiple high-risk application contexts. Their characteristics are summarized in Table .

Table  summarizes their key properties,
highlighting the demographic imbalances and class skew that motivate our
fairness-preserving architecture.

These datasets exhibit varying degrees of class imbalance (24-45% positive class) and demographic skew. The NIJ and COMPAS datasets originate from actual correctional risk assessment systems, making them particularly relevant for validating our architecture's real-world applicability. Adult Census and German Credit provide cross-domain validation in financial decision-making contexts, where similar fairness concerns arise. Figure  illustrates the demographic distributions and imbalances across the four benchmark datasets.

% Experimental Setup

Experiments were executed on a Spark-based cluster configured with four worker nodes (16 vCPUs each). Spark handles the ingestion and Bronze-layer transformations, while the Silver and Gold layer components for causal discovery, DP-SGD synthetic data generation, and privacy verification are implemented as modular Python workflows orchestrated within the Spark-driven pipeline, without deploying the Ray-based execution model that is proposed in the reference architecture for future large-scale deployments.

To align with regulatory standards, the pipeline was configured with strict parameters:

- Privacy Budget: $\epsilon < 1.0$ (GDPR-aligned threshold),
- Minimum anonymity guarantee: $k=5$,
- Anonymization: Adaptive application of $k$-anonymity, $l$-diversity, and $t$-closeness,
- Causal refutation: Performed using DoWhy,
- Bias analysis: Conducted using AIF360,
- Privacy metrics: Computed with ARX.

We evaluate the FAIR-CARE pipeline relative to two baselines representing common industry practices:

- Baseline A (Naive ETL): Standard pseudonymization using hashed identifiers, without statistical anonymization, causal filtering, or fairness enforcement.
- Baseline B (Masking): Static suppression-based masking applied prior to modeling, without synthetic generation or causal validation.

These baselines reflect the typical configurations found in production ETL pipelines and allow us to assess how architectural governance mechanisms influence privacy, fairness, and utility.

% Results and Discussion

% Experimental Outcomes
Our evaluation goes throughout the three pillars of the CARE framework, Privacy, Fairness, and Utility and how it impacts the FAIR-CARE Lakehouse itself. Privacy is quantified through the achieved Differential Privacy budget ($\epsilon$) and $k$-anonymity thresholds, while fairness is assessed using Equalized Odds Difference (EOD) and Demographic Parity Difference (DPD), and utility is measured through AUC. Together, these metrics enables an overall assessment of whether the architecture supports ethically governed AI in high-risk domains.

% Benchmark Comparisons Across Fairness Datasets
The FAIR-CARE pipeline was tested on three fairness datasets: COMPAS, Adult Census, and German Credit. In all cases, the architecture showed an improvement over the baseline in fairness consistency issues. In particular, an ablation study performed for the COMPAS dataset is reported in Table .

The FAIR-CARE configurations proved to consistently outperform the baseline for all experiments, with Config A ($k$-anonymity with $k=5$) and Config B (Differential Privacy with $\epsilon=1.0$) achieving the highest FAIR-CARE Scores of 0.98. It is important to note that the privacy risk increased from 100% (no protection) for the baseline to 6.7% for $k$-anonymity, with FAIR-CARE maintaining utility at 1.00. This shows that privacy-preserving methods can be applied to high-risk data about recidivism without presenting notable accuracy drawback. The causal configuration (Config C) also reported an improvement, validating the importance of correlation discovery to overall governance . In all cases, these findings shows that the FAIR-CARE architecture can be used to provide privacy and fairness guarantees with limited utility loss.

Figure  shows a summary of layer-by-layer FAIR-CARE Scores for all four datasets, demonstrating the FAIR-CARE architecture to be able to improve Bronze, Silver, and Gold governance levels with the overall FAIR-CARE Score in a state of high-readiness.

The impact of these configurations on the composite FAIR-CARE Score is visualized in Figure . While all configurations maintain high utility, the primary benefit of Configs A and B is the substantial reduction in privacy risk, demonstrating that ethical safeguards can be implemented without compromising model performance.

As previously shown in Figure , the demographic imbalances are present in the benchmark datasets. The gender skew (67--88% male across datasets) and racial composition variations (from 51% African-American in COMPAS to 85% White in Adult) demonstrate why traditional ETL pipelines fail: models trained on these distributions will inherently encode historical biases . The FAIR-CARE architecture addresses this through three mechanisms: (1) PII detection and pseudonymization in the Bronze Layer, (2) causal filtering to remove demographic proxies in the Silver Layer, and (3) reweighing and bias mitigation in the Gold Layer.

% Causal Validity Across Pipelines
As part of the Silver Layer processing, we examined how causal filtering behaves using DoWhy refutation tests. During causal discovery, we were explicitly testing for the direction of influence between sensitive attributes and outcomes. This allows the pipeline to flag any potential false correlations, ZIP codes standing in for race, that most typical ETL workflows would swallow uncritically. This addition of causal validation strengthens the overall FAIR-CARE score, since models that lean on unverified assumptions are penalized. The inclusion of causal validation contributes meaningfully to the composite FAIR-CARE score by penalizing models that rely on unverified assumptions.

% Case Study: The NIJ Recidivism Challenge
Using the causal diagram in Figure \ref {fig:nij_causal_dag}, we trace how sensitive attributes, supervision practices, and outcomes relate to each other in the NIJ Recidivism Forecasting Challenge dataset. The highlighted path from race to recidivism, mediated by the level of supervision, suggests one pathway by which differences in supervision intensity, rather than actual risk, could create a structural bias in our observations.

In an attempt to recognize this in a real high-stakes setting, we test Logistic Regression on the NIJ Recidivism Forecasting Challenge dataset . We use ROC AUC as the primary usefulness metric for measuring performance while tracking the fairness metrics Statistical Parity Difference and Equalized Odds Difference. As mentioned earlier, the organizers emphasize the goal of limiting error-rate disparities across racial groups, pointing to the disproportionate harm caused by false positives in decisions to supervise or sanction.

- Baseline (Legacy): Direct modeling on raw data achieved a high AUC of 0.92 but exhibited substantial privacy leakage and large group disparities (e.g., SPD and EOD both around 0.18), indicating strong predictive signal at the cost of fairness.
- Silver Layer (Privacy): Applying the Privacy Preservation module with $\epsilon = 1.0$ (a strict, regulation-aligned privacy budget) maintained a robust AUC of 0.91, showing that most predictive utility can be preserved under differential privacy.
- Gold Layer (Fairness): After bias mitigation via AIF360 Reweighing, the model achieved a marked reduction in fairness disparity (SPD and EOD reduced to approximately 0.07) with only a minor decrease in AUC (final AUC $\approx 0.91$), demonstrating that the architecture can attenuate historical bias while largely retaining predictive performance.

Although the FAIR-CARE pipeline does not maximize predictive accuracy, it yields a more ethically robust model that explicitly trades a small amount of utility for stronger protections against discriminatory and privacy-violating behavior, consistent with emerging principles for trustworthy, high-risk AI systems. This trade-off is summarized in Figure , which visualizes the joint evolution of AUC and fairness metrics across the Baseline, Silver, and Gold configurations.

% FAIR-CARE Score Analysis
The FAIR-CARE Score integrates ethical and technical metrics across the Medallion layers. As seen in Table , all COMPAS configurations have FAIR-CARE Scores above 0.97, indicating that the naive baseline already shows some level of technical utility but lacks clear privacy and causal protection. Adding FAIR-CARE configurations (Configs A--C) increases the overall score to about 0.98 while adding architectural constraints for privacy and causal validity.

The baseline is given a FAIR-CARE Score of 0.97, largely due to its high predictive accuracy and lack of serious failure on governance tasks, but with a privacy risk of 100% due to the absence of anonymization or differential privacy mechanisms. In contrast, the privacy-enhanced configurations lower the privacy risk considerably, to 6.7% for $k$-anonymity and 10% for differential privacy, all while keeping a Silver Score of 1.00 and preserving downstream utility at 1.00. This illustrates that ethical readiness is not an accidental outcome of individualized privacy or fairness methods, but rather a product of combined architectural constraints across the Bronze, Silver, and Gold layers.

For ease of interpretation, a score above 0.95 can be seen as reflecting high levels of readiness for governance on this metric, independent of the specific configuration. On this qualitative definition, all COMPAS variants fall within the same high-readiness region, but the FAIR-CARE configurations provide a more optimal trade-off between near-maximal scores for utility and lower privacy risk of exposure. On this interpretation, the FAIR-CARE architecture transforms the technically strong but weakly governed baseline to one that maintains strong levels of performance while enhancing privacy, fairness, and causal validity guarantees.

Figure  compares FAIR-CARE Scores and utility--privacy trade-offs across four privacy techniques, showing that Differential Privacy achieves the highest composite score and utility retention but at the cost of higher residual privacy risk, while classical anonymization methods offer marginally lower FAIR-CARE Scores with similar utility and lower estimated privacy exposure.

% Legal and Ethical Compliance
We now evaluate the architectural implications in the context of regulatory frameworks, addressing RQ3:

- GDPR: Differential Privacy ($\epsilon \le 1.0$) directly supports the GDPR definition of "anonymization" by preventing singling out, linkability, and inference .
- HIPAA: The Silver Layer implements the ``Expert Determination'' pathway through DP mechanisms and statistical disclosure control, ensuring that residual re-identification risk is ``very small'' .
- CCPA/CPRA: Functional separation between the Bronze (identifier keys) and Gold (analytic features) layers ensures that downstream consumers cannot reconstruct identity, satisfying statutory separation requirements.

Collectively, these results demonstrate that the FAIR-CARE Lakehouse not only improves fairness and privacy performance but also provides a principled architectural mechanism for satisfying regulatory mandates in high-risk AI systems.

% Limitations
While progress has been made, many gaps are still apparent. Construction of the Causal Graph in the Silver Layer depends on input from domain experts. We do make use of algorithmic checks to help, but fully automated causal discovery is an open research challenge. Our current setup focuses on validation rather than automatic feature suppression. As the experiment by NIJ illustrates, strict enforcement of fairness constraints can lower performance. Where high accuracy is essential, as in violent crime predictions, this type of trade-off requires thoughtful policy decisions well beyond the realm of software design. Our evaluation so far has been on structured tabular data. Processing unstructured modalities, such as video or audio, would require different privacy measures-for instance, deepfakes for anonymization-which are beyond the scope of this study.

% Artifact Availability
To support reproducibility, we provide an anonymized code repository with the full FAIR-CARE pipeline implementation and experiment scripts at the following URL (anonymized for review): https://anonymous.4open.science/r/fair-care-data-pipeline-ADDA.

% Threats to Validity

Despite the structured design of the FAIR-CARE Lakehouse and its alignment with regulatory, causal, and ethical principles, several threats to validity must be acknowledged. These factors affect the interpretation, generalizability, and robustness of the study's findings.

% Internal Validity
Internal validity concerns whether the observed improvements can be reliably attributed to the architecture.
A first threat is causal graph correctness: the causal graph construction in the Silver Layer depends on domain expert elicitation and algorithmic discovery (e.g., PC, NOTEARS), and these algorithms are sensitive to hyperparameter choices and may be influenced by unobserved confounders, so flawed structures can lead to ineffective or misleading feature filtering in the Gold Layer.
To mitigate this, we relied on refutation tests (e.g., placebo treatments) via DoWhy, but automated causal discovery and validation remain open research challenges.
A second threat is sensitivity to privacy parameters: the Differential Privacy mechanisms (e.g., Laplace noise) introduce stochasticity, so small deviations in the privacy budget ($\epsilon$) or clipping thresholds can yield different utility–fairness trade-offs, and although we fixed random seeds for reproducibility, single-run results may still vary slightly.

% External Validity
External validity pertains to the generalizability of our findings beyond the datasets and contexts studied.
Our evaluation focuses exclusively on structured tabular datasets, which means the behavior of the architecture for unstructured modalities such as video or free-text narratives remains unknown and would require distinct privacy mechanisms (e.g., redaction, blurring) that are not yet implemented in the Silver Layer.
Moreover, the datasets used (COMPAS, NIJ, US Census–derived Adult, German Credit) originate from U.S. or closely related socio-legal contexts, so bias patterns and regulatory requirements may differ significantly in international domains (e.g., EU healthcare or social services), and the specific causal structures and fairness–utility trade-offs observed here may not transfer directly.

% Inference Validity
Construct validity examines whether our measures accurately reflect the ethical constructs under study. We assessed fairness using group-based measures, particularly the population parity difference (DPD) and the equal opportunity difference (EOD), which are standard, but individual fairness does not capture procedural fairness or long-term social impact. Similarly, we operationalized utility primarily as AUC, which does not reflect calibration, interpretability, or the influence of downstream decision-making that may be essential in high-risk domains. Furthermore, the FAIR-CARE score is a composite measure whose weighting scheme ($w_S, w_G$) is architecturally defined. Although this weighting is meaningful for comparing configurations in our study, alternative weightings may be needed given different stakeholder preferences or regulatory interpretations, which could lead to different conclusions about the ethical fitness of a particular configuration.

% Validity of Conclusions
Conclusion validity relates to the statistical strength of our claims. Our quantitative evaluation is based on a limited number of benchmark datasets ($N=4$), and although we observed consistent trends, some paired comparisons lacked statistical significance at the conventional $p < 0.05$ threshold, limiting the strength of inferential claims. Furthermore, the fairness gains demonstrated on static benchmarks may differ in production environments subject to distributional shift (data drift), irregular update cycles, or adversarial behavior. As a result, the real-world robustness of the FAIR-CARE Lakehouse remains to be validated in longitudinal deployments, and future work should incorporate more extensive repetitions, larger and more diverse datasets, and robustness checks under realistic operational conditions.

% Conclusion and Future Work
This paper presented the FAIR-CARE Lakehouse, a reference architecture that unifies FAIR data management principles with ethical governance mechanisms tailored for high-risk AI systems. By extending the Medallion model with privacy-preserving technologies, causal inference components, and lineage-aware enforcement policies, the architecture reframes the data pipeline as a socio-technical system in which fairness, privacy, and regulatory compliance are first-class architectural concerns. In addressing our research questions, we demonstrated that ethical governance can be enforced end-to-end through architectural mechanisms embedded in each stage of the pipeline (RQ1). The combination of differential privacy, synthetic data generation, causal DAG validation, and fairness-aware feature selection forms a technically grounded toolkit that enables enforcement of CARE constraints while maintaining competitive utility (RQ2). Furthermore, we introduced the FAIR-CARE score as a composite metric that operationalizes ethical readiness and provides a measurable basis for evaluating datasets in high-risk domains (RQ3). The contributions noted above suggest that architectural design, rather than downstream model tuning, is a good place to embed principles of trustworthy AI. Overall, FAIR-CARE Lakehouse advances the role of software architecture in AI governance by outlining a structured, actionable, and measurable approach to designing an ethical data pipeline. This architecture provides a basis for building ongoing governance mechanisms and future adaptability, and supports the long-term evolution of responsible AI infrastructures.

% Future Work

The key step in the architecture of FAIR-CARE is embracing FL to enable distributed model building with privacy across organizations. Correction data lives in various jurisdictions, each having its own legal, operational, and infrastructure limits. Due to trust between organizations and regulatory rules, pulling all this data into a single place is usually not possible; for example, GDPR has rules on data minimization. FL frameworks are the way collaborators can train models together without exchanging raw data, hence enhancing the system's privacy guarantees: NVFlare, Flower, or Ray Fed. From an architectural perspective, federated learning creates new frontiers for governance: privacy budgets must be managed across federated rounds, fairness must be monitored across participating nodes, and causal structures differ between silos. In the future, it will be particularly useful to generalize the FAIR-CARE score to a federated governance score capturing heterogeneity and deviations from ethical conformance across sites. Automated governance tools, such as node-level anomaly detection, fair deviation monitoring, and DP-aware aggregation, will be crucial enablers of ethical conformance in asynchronous competitive training regimes. Thus, integrating FL and Lakehouse FAIR-CARE is not simply a technical upgrade but rather one aspect of an intrinsically evolving architecture for high-risk AI ecosystems, where data sources evolve independently and governance has to adapt in real time.
